% Teaching manuscript. Edit this TeX source; Markdown is compatibility navigation only.

\StudyChapter{CS5489-Lecture10}{Lecture 10 · Autoencoders \& Deep Generative Models}
\StudyTarget{CS5489-Lecture10__lecture-10-autoencoders-generative}


\StudyHeading{1}{从压缩重构到无中生有：AE、VAE、GAN 与扩散模型}
\StudyTarget{CS5489-Lecture10__intro}

\StudyRepoRef{CS5489/course-notes/README.md}{}{课程目录} · \StudyRef{CS5489-Lecture09}{root}{上一讲：聚类分析} · \StudyRef{CS5489-Lecture11}{root}{下一讲：大语言模型}

前面两讲用 PCA、KPCA 做降维，用 K-Means、GMM 做聚类。那能不能直接让神经网络在没有标签的情况下学会压缩特征、去噪修复，甚至“照着数据分布凭空画出逼真的新图片”？

这一讲沿着“表征学习\ensuremath{\to}概率生成”的阶梯一路攀登：前半讲（Lecture10a）从最基础的\textbf{自编码器（Autoencoder, AE）}讲起，解决恒等退化的\textbf{去噪自编码器（DAE）}、面向图像的\textbf{卷积自编码器（ConvAE）}与 \textbf{Masked Autoencoder（MAE）}，以及结构同源的语义分割网络 \textbf{FCN 与 U-Net}；后半讲（Lecture10b \& 10c）进入\textbf{深度生成模型}的核心——先用\textbf{重参数化技巧（Reparameterization Trick）}打通随机采样的反向传播，再依次拆解\textbf{变分自编码器（VAE）}、\textbf{生成对抗网络（GAN）}与当今文生图霸主\textbf{扩散模型（Diffusion Models / DDPM / DDIM / CFG）}。

\textbf{基础复习（可跳过）：} 如果对 CNN 卷积与池化不熟，先看\StudyRef{CS5489-Lecture06}{cnn-layers}{Lecture 6 卷积层}；如对 PCA 重构目标不熟，可对照\StudyRef{CS5489-Lecture08}{pca}{Lecture 8 PCA}。

学习目标与难点见\StudyRef{CS5489-Lecture10}{learning-priorities}{文末学习要求表}。


\StudyHeading{1}{Exam focus｜这一讲怎样安排复习}
\StudyTarget{CS5489-Lecture10__exam-review}
\StudyTarget{CS5489-Lecture10__exam-focus}

\textbf{本讲在已辨识的 3 套历史期末卷（2020B、2021A、2021B）中每套卷都占了整整 15 分（1 道多选 + 1 道 10 分大题）！} AE、VAE、GAN 全部是高频大题与多选题主力；近年新增的 MAE、U-Net 与 Diffusion（前向／反向公式、噪声预测损失、CFG 引导）也在 2026 课件里占了极大篇幅。

\begingroup\small\setlength{\tabcolsep}{3pt}\renewcommand{\arraystretch}{1.22}
\begin{longtable}{@{}>{\raggedright\arraybackslash}p{\dimexpr0.23\linewidth-1.50000\tabcolsep\relax}>{\raggedright\arraybackslash}p{\dimexpr0.35\linewidth-1.50000\tabcolsep\relax}>{\raggedright\arraybackslash}p{\dimexpr0.22\linewidth-1.50000\tabcolsep\relax}>{\raggedright\arraybackslash}p{\dimexpr0.20\linewidth-1.50000\tabcolsep\relax}@{}}
\toprule
考点 & 复习动作 & 历史期末（共3套） & 入口 \\
\midrule\endfirsthead
考点 & 复习动作 & 历史期末（共3套） & 入口 \\
\midrule\endhead
\bottomrule\endfoot
Autoencoder、DAE、MAE 与 U-Net & 记牢 AE 无监督、不必镜像对称、比 PCA 慢；可用预训练层；DAE 抗恒等退化 & 直接 1 套 & \StudyRef{CS5489-Lecture10}{autoencoders}{讲解} \\
重参数化技巧与似然不可算原因 & 记住 $y=\mu+\sigma z$ 拆出随机性；变量替换需 $f^{-1}$ 可逆且可微 & 直接 2 套 & \StudyRef{CS5489-Lecture10}{reparam}{讲解} \\
VAE 损失、与 AE 异同及单类分类 & 对比 AE vs VAE（后验分布、KL 正则、采样插值）；用 VAE 隐空间做苹果单类异常检测 & 直接 3 套 & \StudyRef{CS5489-Lecture10}{vae}{讲解} \\
GAN 目标、非饱和损失与模式崩塌 & 记牢 $D^*=\frac{p_d}{p_d+p_m}$、$-\log D(G(z))$ 梯度更稳、模式崩塌定义及无显式似然 & 直接 2 套 & \StudyRef{CS5489-Lecture10}{gan}{讲解} \\
扩散模型（DDPM / Latent / CFG） & 记住前向闭式 $x_t$、U-Net 预测噪声 $\mathbb{E}\|\epsilon-\epsilon_\theta\|^2$ 与 CFG 线性外推 & 近年新增模块 & \StudyRef{CS5489-Lecture10}{diffusion}{讲解} \\
\end{longtable}\endgroup

口径说明：统计已辨识的 3 套独立期末卷（2020B、2021A、2021B），同卷不重复计数；扩散模型为新版课件重点扩充内容。完整题号见\StudyRef{CS5489-Lecture10}{exam-topic-index}{讲末历史题源表}。


\StudyHeading{1}{1. 自编码器家族：AE、DAE、ConvAE、MAE 与 U-Net}
\StudyTarget{CS5489-Lecture10__autoencoders}

考点：Autoencoder 性质与去噪机制 \StudyBadge{final}{期末直接 · 1 套}

\textbf{自编码器（Autoencoder, AE）}由编码器 $\mathbf{z}=f(\mathbf{x})$ 和解码器 $\hat{\mathbf{x}}=g(\mathbf{z})$ 串联而成：中间较窄的隐藏层 $\mathbf{z}$ 叫\textbf{瓶颈层（bottleneck / code）}，强迫网络把高维输入先压进低维隐空间，再尽量把原始输入还原出来。它完全\textbf{不需要类别标签}（目标就是输入 $\mathbf{x}$ 本身），按输入范围选择重构损失：实数输入用均方误差 $\|\mathbf{x}-\hat{\mathbf{x}}\|_2^2$，落在 $[0,1]$ 的灰度像素常用二元交叉熵 $-\sum\bigl[x_j\log\hat{x}_j+(1-x_j)\log(1-\hat{x}_j)\bigr]$。

对照 2021A 期末 Q6，以下五条性质务必分清正误：
\begin{itemize}
\item \textbf{与 PCA 的关系及效率：} 若单层 AE 的编解码采用转置权重共享 $\mathbf{z}=\mathbf{A}^T\mathbf{x},\;\hat{\mathbf{x}}=\mathbf{A}\mathbf{z}$ 且用线性激活，它就近似于（不含正交约束的）PCA；加多层与非线性后能学非线性流形。但\textbf{训练 AE 比算 PCA 慢得多}（PCA 有闭式 SVD 解，AE 需要多轮迭代反向传播，故 2021A Q6 选项 A“学 AE 比学 PCA 更高效”为错）。
\item \textbf{主要用途：} 降维（取编码器输出 $\mathbf{z}$）与图像去噪（选项 B 正确）。
\item \textbf{完全无监督／自监督：} 训练只需原始输入 $\mathbf{x}$，不需要部分或全部类别标签（选项 C 为错）。
\item \textbf{支持迁移增强且不要求镜像对称：} 可以把其他任务预训练好的骨干层拿来增强编码器或解码器（选项 D 正确）；并且编码器与解码器\textbf{不必严格镜像对称}（选项 E 为错），两者甚至不必共享权重（2021B Q6 选项 B 同样强调解码器无需与编码器共享权重）。
\end{itemize}

\StudyHeading{2}{去噪自编码器（DAE）、卷积 AE、MAE 与语义分割 U-Net}
\StudyTarget{CS5489-Lecture10__dae-mae-unet}

标准 AE 有一个隐患：如果瓶颈层足够宽或网络容量太大，它可以轻而易举学会一个毫无用处的\textbf{恒等函数（identity function）} $g(f(\mathbf{x}))=\mathbf{x}$，原封不动把像素搬过去，根本没学到数据分布结构！两种常见解法：(1) 给隐码加稀疏正则 $L=\|\mathbf{x}-g(f(\mathbf{x}))\|_2^2 + \lambda\|f(\mathbf{x})\|_1$；(2) \textbf{去噪自编码器（Denoising Autoencoder, DAE）}——先故意把输入破坏成 $\tilde{\mathbf{x}}$（如加高斯噪声，或直接对输入层用 Dropout 随机清零部分像素），再让网络重构\emph{干净}原图 $\mathbf{x}$：
\[
\min_{f,g}\;\sum_i \|\mathbf{x}_i - g(f(\tilde{\mathbf{x}}_i))\|_2^2.
\tag{10.1}
\]
这样网络即使隐码维度比输入还大（过完备）也无法靠背恒等映射偷懒，必须学会把偏离数据流形的破损点 $\tilde{\mathbf{x}}$ 拉回干净流形 $\mathbf{x}$ 上（几何上是学习指向数据流形的\textbf{向量场}，概率上是最大化 $\mathbb{E}_{\mathbf{x},\tilde{\mathbf{x}}}[\log p_{\text{decoder}}(\mathbf{x}\mid f(\tilde{\mathbf{x}}))]$）。注意测试提取干净图像特征时，输入 Dropout 关闭不加噪。

把全连接层换成卷积和池化（解码端用上采样恢复尺寸），就得到\textbf{卷积自编码器（ConvAE）}，在隐空间两点间做线性插值 $\hat{\mathbf{z}}=(1-\alpha)\mathbf{z}_1+\alpha\mathbf{z}_2$ 再解码，就能看到数字“7”平滑变形为“1”。进一步把这一思想搬到 Vision Transformer 上，就诞生了自监督里程碑 \textbf{Masked Autoencoder（MAE）}：随机遮住 \textbf{75\%} 的图像块（patches），只让 ViT 编码器看剩下的 25\% 可见块，再由解码器把所有被遮盖的块像素重建出来；在 ImageNet-1k 无标签预训练后微调，效果甚至超过直接在 ImageNet-1k 上有监督训练的 ViT（但仍逊于超大标注集 JFT 预训练）。

结构上，\textbf{编码器–解码器}也直接催生了像素级分类任务——\textbf{语义分割（Semantic Segmentation）}（给图像每个像素打类别标签）：\textbf{全卷积网络（FCN）}去掉末端 MLP、先下采样提语义再上采样出分割图；\textbf{U-Net} 在 FCN 基础上加入\textbf{跨层跳跃拼接（skip concatenation）}，把编码端高分辨率浅层特征直接拼到解码端对应上采样层，同时兼顾高层语义与精细物体边界定位（这一 U-Net 骨干后来也成了扩散模型的核心去噪网络！）。

来源：Lecture10a，第 1–98 个单元（AE、权重共享、DAE 两种解释、ConvAE 隐空间插值、MAE、FCN 与 U-Net）。


\StudyHeading{1}{2. 重参数化技巧与深度生成模型基石}
\StudyTarget{CS5489-Lecture10__reparam}

考点：重参数化技巧与概率密度不可算原因 \StudyBadge{final}{期末直接 · 2 套}

判别分类器干的是“多对一”映射（无数张手写“2”都映射到同一个标签 2）；而\textbf{生成模型（generative model）}要干的是“一对多”映射（给一个条件甚至只给随机种子，就能生成千变万化但都合理的“2”）。普通确定性神经网络遇到一对多目标时，为了最小化平方误差只能输出所有可能结果的平均——比如把不同角度的“2”叠成一团模糊鬼影。

于是我们让神经网络输出分布参数，比如高斯均值与标准差 $y\sim\mathcal{N}(\mu(\mathbf{x}),\sigma(\mathbf{x})^2)$。可新难题来了：计算图里一旦出现“从 $\mathcal{N}(\mu,\sigma^2)$ 随机采样”，梯度没法穿过随机采样节点反传给 $\mu$ 和 $\sigma$！\textbf{重参数化技巧（Reparameterization Trick）}用一个极其优雅的改写彻底解决：把随机性剥离出来交给一个\textbf{参数固定的外部标准噪声} $z\sim\mathcal{N}(0,1)$，再通过确定性可微函数把噪声变成目标样本：
\[
y = \mu(\mathbf{x}) + \sigma(\mathbf{x})\,z, \qquad z \sim \mathcal{N}(0, 1).
\tag{10.2}
\]
现在 $z$ 与网络参数无关，只是一个额外输入；反向传播关于 $\mu(\mathbf{x})$ 和 $\sigma(\mathbf{x})$ 的偏导数一目了然（分别为 $1$ 和 $z$）！推广到一般复杂分布，我们可以用任意神经网络 $f$ 把简单噪声 $z\sim p_z(z)$（例如 $\mathcal{N}(0,\mathbf{I})$ 或均匀分布 $U(0,1)$）变换成复杂目标分布的样本 $y = f(z)$。

\textbf{高频 10 分大题（2021B 期末 Q13）：} 既然可以用神经网络 $y = f(z),\; z\sim\mathcal{N}(0,\mathbf{I})$ 轻松采样复杂分布，那为什么我们通常\textbf{算不出生成分布的精确概率密度（似然函数）$p_y(y)$}？需要什么条件才能算出来？
由概率论的变量替换公式（概率质量守恒 $|p_y(y)\,dy| = |p_z(z)\,dz|$），一维（及多维雅可比推广）下密度满足：
\[
p_y(y) = \left|\frac{df^{-1}(y)}{dy}\right|\,p_z\bigl(f^{-1}(y)\bigr).
\tag{10.3}
\]
（特例验算：若 $z\sim\mathcal{N}(0,1)$ 且 $y=2z-1$，反函数 $z=(y+1)/2$、导数 $1/2$，代入式 10.3 恰好得 $\frac{1}{2}\mathcal{N}((y+1)/2\mid 0,1)=\mathcal{N}(y\mid -1,4)$；若 $z\sim U(0,1)$，学习 $f$ 就等价于学习 $y$ 的逆累积分布函数 $F_y^{-1}$！）
\begin{itemize}
\item \textbf{为什么一般算不出来：} 式 10.3 需要求神经网络 $f$ 的\textbf{逆映射 $f^{-1}$}（以及逆映射的雅可比行列式），而一般的神经网络（含降维、多对一映射、ReLU 等）\textbf{并不可逆（not invertible）}！
\item \textbf{使其可算的条件：} 神经网络 $f$ 必须完全由\textbf{可逆（invertible）且可微（differentiable）}的变换复合而成（这正是标准化流 Normalizing Flows 的基本思想）。
\end{itemize}

来源：Lecture10b，第 12–21 个单元（一对多生成动机、重参数化技巧、变量替换密度公式与逆 CDF 解释）。


\StudyHeading{1}{3. 变分自编码器（VAE）｜把隐空间变成连续平滑的概率地图}
\StudyTarget{CS5489-Lecture10__vae}

考点：VAE 原理、与 AE 对比及单类异常检测 \StudyBadge{final}{期末直接 · 3 套}

普通 AE 只是把每张训练图硬映射到隐空间里的一个孤立点 $\mathbf{z}_i$；点与点之间的空白区域解码器在训练时从未见过，一旦随机抽个隐向量或在两张图之间插值，解码器很容易吐出支离破碎的怪图。

\textbf{变分自编码器（Variational AutoEncoder, VAE）}把隐变量 $\mathbf{z}$ 变成随机变量：设定简单先验 $p(\mathbf{z})=\mathcal{N}(\mathbf{0},\mathbf{I})$，让编码器网络不再输出单个点，而是输出后验高斯分布的参数——均值向量 $\boldsymbol{\mu}(\mathbf{x}_i)$ 与对数标准差 $\log\boldsymbol{\sigma}(\mathbf{x}_i)$，构成近似后验 $q(\mathbf{z}\mid\mathbf{x}_i)=\mathcal{N}\bigl(\boldsymbol{\mu}(\mathbf{x}_i),\,\mathrm{diag}(\boldsymbol{\sigma}(\mathbf{x}_i)^2)\bigr)$。训练时用重参数化技巧采样 $\mathbf{z}_{il}=\boldsymbol{\mu}(\mathbf{x}_i)+\boldsymbol{\sigma}(\mathbf{x}_i)\odot\boldsymbol{\epsilon}_{il}$（$\boldsymbol{\epsilon}_{il}\sim\mathcal{N}(\mathbf{0},\mathbf{I})$）送入解码器（图 10.1 左）。由于精确边际似然 $\log p(\mathbf{x})=\log\int p(\mathbf{x}\mid\mathbf{z})p(\mathbf{z})\,d\mathbf{z}$ 无法直接积分，我们最大化变分下界（ELBO），等价于最小化如下两项之和的 \textbf{VAE 损失函数}：
\[
\mathcal{L}_{\text{VAE}}(\mathbf{x}_i) = \underbrace{-\frac{1}{L}\sum_{l=1}^L \log p(\mathbf{x}_i \mid \mathbf{z}_{il})}_{\text{重构损失（如 MSE / BCE）}}
+ \underbrace{\mathrm{KL}\bigl(q(\mathbf{z}\mid\mathbf{x}_i)\,\|\,p(\mathbf{z})\bigr)}_{\text{正则项：拉近先验 } \mathcal{N}(\mathbf{0},\mathbf{I})}.
\tag{10.4}
\]
其中高斯后验与标准正态先验之间的 KL 散度有简洁闭式解：对单维 $(\mu,\sigma)$ 为 $\frac{1}{2}(\mu^2 + \sigma^2 - 1 - \log\sigma^2)$（例如 $\mu=0,\sigma=1$ 时 KL 为 $0$；$\mu=1,\sigma=2$ 时 KL 为 $2-\log 2\approx 1.3069$）。注意编码器与解码器是\textbf{联合端到端训练}的（并非解耦分开训，2021B Q6 选项 D 为错），也不需要类别标签（选项 E 为错），且解码器不必和编码器共享权重（选项 B 正确）。

\textbf{期末必背对照（2020B 期末 Q12，10 分大题）：AE 与 VAE 的异同及意义}
\begin{itemize}
\item \textbf{相同点：} 两者都把输入编码进低维隐空间再解码重构原输入；都使用重构损失（MSE 或二元交叉熵）训练，且均属无监督学习。
\item \textbf{不同点：} (1) AE 输出确定性的单个隐向量 $\mathbf{z}$，VAE 输出隐向量 $\mathbf{z}$ 的\textbf{高斯后验分布}（均值与方差）；(2) VAE 损失额外包含 \textbf{KL 正则项}，把各样本的后验拉向标准高斯先验 $\mathcal{N}(\mathbf{0},\mathbf{I})$；(3) 训练时 VAE 每次从后验\textbf{随机采样}（给隐向量加噪）后再送进解码器。
\item \textbf{为什么这些差异至关重要：} 因为训练时不断在 $\boldsymbol{\mu}(\mathbf{x}_i)$ 周围随机采样，解码器见过整个连续邻域里的各种隐向量，防止对单个孤立编码过拟合，更鲁棒；加上 KL 项把所有样本后验紧密组织在标准高斯球周围，消除了“空洞”，使隐空间\textbf{连续平滑、非常适合样本间插值与直接从 $\mathcal{N}(\mathbf{0},\mathbf{I})$ 采样生成全新样本}。
\end{itemize}

\textbf{实战迁移（2021A 期末 Q10，10 分大题）：} 手头只有大量\emph{苹果}图片、没有任何\emph{非苹果}图片，怎么用 VAE 判断新图是不是苹果？
思路非常清晰——做成\textbf{单类分类（one-class classification）}：用所有苹果图片训练好 VAE，用编码器算出全部训练苹果图的隐编码 $\{\mathbf{z}_i\}_{i=1}^N$（即后验均值 $\boldsymbol{\mu}(\mathbf{x}_i)$）。推理时，把测试图送进编码器得隐编码 $\mathbf{z}^*$，计算它到所有苹果训练编码的最小距离 $\min_i \|\mathbf{z}^*-\mathbf{z}_i\|$（实践上也可结合重构误差作为辅助信号）：若小于设定阈值则判为“苹果”，否则判为“非苹果”。

\StudyImage{CS5489/course-notes/assets/l10-generative-models.png}
\StudyCaption{图 10.1：左图展示 VAE 重参数化计算图——将随机采样 $\boldsymbol{\epsilon}\sim\mathcal{N}(\mathbf{0},\mathbf{I})$ 变为外部输入，使梯度能畅通传回编码器 $\boldsymbol{\mu},\boldsymbol{\sigma}$；右图对比 GAN 原始极小极大损失与非饱和损失对判别器 logit 的梯度幅值，早期假样本被一眼识破（$D(G(z))\approx 0$）时原始损失梯度近乎消失，非饱和损失梯度仍强劲（教学自绘）。}

来源：Lecture10b，第 22–83 个单元（VAE 原理、ELBO 损失、MLP-VAE 与 Conv-VAE 在 MNIST 上的隐空间插值）。


\StudyHeading{1}{4. 生成对抗网络（GAN）｜造假者与鉴宝师的极小极大博弈}
\StudyTarget{CS5489-Lecture10__gan}

考点：GAN 博弈目标、非饱和损失与模式崩塌 \StudyBadge{final}{期末直接 · 2 套}

既然用重参数化 $x=G(z),\; z\sim\mathcal{N}(0,\mathbf{I})$ 无法直接写出似然 $p_m(x)$，\textbf{生成对抗网络（Generative Adversarial Networks, GAN）}干脆绕开显式密度估计，直接看生成样本像不像真货！它同时训练两个互相博弈的网络：
\begin{itemize}
\item \textbf{生成器（Generator, $G$）：} 吃进随机噪声 $z\sim\mathcal{N}(0,\mathbf{I})$，造出假样本 $G(z)$，目标是骗过判别器；
\item \textbf{判别器（Discriminator, $D$）：} 二分类器，输出 $D(x)=p(\text{real}\mid x)\in(0,1)$，尽量给真样本打 1、给假样本打 0。
\end{itemize}
两者构成经典的\textbf{极小极大零和博弈（minimax game）}：
\[
\min_{G}\max_{D}\; V(D,G) = \mathbb{E}_{x\sim p_d}[\log D(x)] + \mathbb{E}_{z\sim\mathcal{N}(0,\mathbf{I})}\bigl[\log\bigl(1 - D(G(z))\bigr)\bigr].
\tag{10.5}
\]
理论分析表明，在容量足够时，固定 $G$ 的最优判别器是 $D^*(x) = \frac{p_d(x)}{p_d(x)+p_m(x)}$；此时最优生成器恰好满足 $p_m^*(x)=p_d(x)$，真假分布完全重合，判别器彻底猜不出来（$D^*(x)=0.5$，双方准确率均停在 50\% 随机猜测，平衡点是一个\textbf{鞍点 saddle point}）。

三条\textbf{期末必背核心考点}（2020B 期末 Q6 全部考到）：
\begin{enumerate}
\item \textbf{GAN 的优缺点：} 优点是不需要近似似然函数或变分推断，也\textbf{不需要计算配分函数（partition function）}，采样极快、图像清晰锐利；缺点是\textbf{无法直接计算样本的概率密度（似然值）}，且鞍点博弈训练极不稳定。
\item \textbf{为什么实践改用非饱和损失 $\min_G -\mathbb{E}_z[\log D(G(z))]$？（图 10.1 右）：} 训练刚开始时，$G$ 造的假货太拙劣，判别器轻轻松松打出 $D(G(z))\approx 0.01$。如果沿用原始极小极大目标 $\log(1-D(G(z)))$，因为 $D=\sigma(u)$ 的导数是 $D(1-D)$，对 logit $u$ 的梯度幅值仅剩 $D = 0.01$（\textbf{梯度消失}，$G$ 几乎学不动）；换成让 $G$ 最小化 $-\log D(G(z))$（等价于给假样本贴上“真”标签算交叉熵），对 logit $u$ 的梯度幅值高达 $1-D = 0.99$（大近 100 倍！），训练稳定性远胜原始目标。
\item \textbf{什么是模式崩塌（Mode Collapse）？} 当真实分布有多个山峰（多模态，如 10 种手写数字）时，生成器可能“偷懒”只学会生成其中\textbf{某一个或少数几个模态}，因为靠这一个逼真的模态就已经能骗过判别器（等判别器识破后，$G$ 又跳去另一个单峰来回打转）。注意：模式崩塌是\textbf{缺乏生成多样性（只覆盖部分模态）}，并\emph{不是}“生成器骗不过判别器”（2020B Q6 选项 E 故意把概念反写设坑！）。
\end{enumerate}

来源：Lecture10b，第 84–193 个单元（GAN 动机、优缺点、极小极大分析、1D 仿真、模式崩塌、DCGAN、CycleGAN 与人脸合成）。


\StudyHeading{1}{5. 扩散模型（Diffusion Models）｜一步步加噪，再一步步请 U-Net 去噪}
\StudyTarget{CS5489-Lecture10__diffusion}

GAN 虽然清晰但难训、易模式崩塌；VAE 好训但单步解码常略显模糊。\textbf{扩散模型（Diffusion Models / DDPM）}换了一条稳健得多的路：把“从纯噪声直接画出一张高清图”拆解成成百上千个极小的微调步骤，每一步都有明确的监督信号！

\textbf{前向加噪过程（Forward diffusion，固定无需学习）：} 给定预设方差表 $\{\beta_1,\dots,\beta_T\}$（记 $\alpha_t=1-\beta_t,\;\bar{\alpha}_t=\prod_{s=1}^t \alpha_s$），每一步把图像淡化一点并混入标准高斯噪声 $\boldsymbol{\epsilon}_t\sim\mathcal{N}(\mathbf{0},\mathbf{I})$：
\[
\mathbf{x}_t = \sqrt{1-\beta_t}\,\mathbf{x}_{t-1} + \sqrt{\beta_t}\,\boldsymbol{\epsilon}_t
\;\Longrightarrow\;
\mathbf{x}_t = \sqrt{\bar{\alpha}_t}\,\mathbf{x}_0 + \sqrt{1-\bar{\alpha}_t}\,\boldsymbol{\epsilon}.
\tag{10.6}
\]
右边的闭式边缘公式（Lecture10c 第 9 个单元）极其好用：训练时不用真的迭代 $t$ 步，只需直接按 $\bar{\alpha}_t$ 一步合成任意时刻 $t$ 的加噪图 $\mathbf{x}_t$（例如两步 $\beta_1=\beta_2=0.1$ 时，$\bar{\alpha}_2=0.81$，信号系数 $\sqrt{0.81}=0.9$，噪声方差 $1-0.81=0.19$）。当 $T$ 足够大时，$\mathbf{x}_T$ 彻底变成纯高斯白噪声。

\textbf{反向去噪与简化训练目标（Reverse process）：} 真实逆转移 $q(\mathbf{x}_{t-1}\mid\mathbf{x}_t)$ 不可直接求，我们用神经网络（通常是 \textbf{U-Net}，吃进当前噪声图 $\mathbf{x}_t$ 和时间步 $t$）拟合高斯逆转移 $p_\theta(\mathbf{x}_{t-1}\mid\mathbf{x}_t)=\mathcal{N}(\boldsymbol{\mu}_\theta(\mathbf{x}_t,t),\boldsymbol{\Sigma}_\theta(\mathbf{x}_t,t))$。最小化负对数似然的变分上界后，Ho 等人（DDPM）将其化简成一个惊人简洁的\textbf{噪声预测均方误差}：
\[
\mathcal{L}_{\text{simple}}(\theta) = \mathbb{E}_{t,\mathbf{x}_0,\boldsymbol{\epsilon}}\bigl\|\boldsymbol{\epsilon} - \boldsymbol{\epsilon}_\theta(\mathbf{x}_t, t)\bigr\|_2^2.
\tag{10.7}
\]
也就是说，网络只需要看一眼含噪图 $\mathbf{x}_t$ 和时刻 $t$，把刚才掺进去的那份高斯噪声 $\boldsymbol{\epsilon}$ 精准猜出来即可！

结合 Lecture10c 实操实验，记住四个现代扩散系统的关键工程概念：
\begin{itemize}
\item \textbf{像素空间 vs 隐空间扩散（Stable Diffusion）：} 小图如 CIFAR-10（$32\times 32\times 3$）可以直接在 RGB 像素空间扩散；高清图（$512\times 512$）计算太重，Stable Diffusion 先用预训练 VAE 把图片压进 $[1, 4, 64, 64]$ 的紧凑\textbf{隐空间（latent space）}做扩散去噪，最后再由 VAE 解码回高清像素。
\item \textbf{加速采样（DDIM）：} 原始 DDPM 需要 1000 步反向去噪，确定性采样器 DDIM（$\eta=0$）只需 20–30 步就能还原高质量图像。
\item \textbf{无分类器引导（Classifier-Free Guidance, CFG）：} 文生图时同时算一次无条件预测 $\boldsymbol{\epsilon}_{\text{uncond}}$ 和一次带文本提示的预测 $\boldsymbol{\epsilon}_{\text{cond}}$，再按引导尺度 $s$ 线性外推：$\tilde{\boldsymbol{\epsilon}} = \boldsymbol{\epsilon}_{\text{uncond}} + s(\boldsymbol{\epsilon}_{\text{cond}} - \boldsymbol{\epsilon}_{\text{uncond}})$（如 $s=7.5$；当 $s\le 1$ 时相当于不放大条件）。调高 $s$ 让图像更紧扣提示词，但过高会导致过饱和、失真。
\item \textbf{空间条件控制与局部编辑：} GLIGEN 允许附加短语与边界框作为\textbf{软空间条件}引导物体布局；局部重绘（Inpainting）则通过二值掩码（白=重绘、黑=保留）与 \StudyInline{strength}（控制起始加噪深度与执行步数）对原图指定区域做精细修改。
\end{itemize}

\textbf{自查 / Check：} 给定 $\boldsymbol{\epsilon}_{\text{uncond}}=[1.0,-0.5]$、$\boldsymbol{\epsilon}_{\text{cond}}=[1.4,-0.1]$，当 CFG 引导系数 $s=7.5$ 时，引导后的噪声预测 $\tilde{\boldsymbol{\epsilon}}$ 是多少？为什么要先把高清图放进 VAE 隐空间做扩散？ / Given $\boldsymbol{\epsilon}_{\text{uncond}}=[1.0,-0.5]$ and $\boldsymbol{\epsilon}_{\text{cond}}=[1.4,-0.1]$, compute the guided noise prediction at CFG scale $s=7.5$. Why does Stable Diffusion run diffusion in VAE latent space?

\begin{StudySolution}{答案 / Answer}
$\tilde{\boldsymbol{\epsilon}} = [1.0,-0.5] + 7.5\times [0.4, 0.4] = [4.0, 2.5]$。在 $512\times 512\times 3$ 像素空间反复跑 U-Net 计算与显存开销过大；先经 VAE 压成 $4\times 64\times 64$ 隐张量，大幅降低每步扩散成本，生成后再由 VAE 解码回像素图。 / $\tilde{\boldsymbol{\epsilon}} = [1.0,-0.5] + 7.5[0.4,0.4] = [4.0,2.5]$. Running iterative U-Net denoising directly on $512\times 512$ pixels is costly; compressing to a $4\times 64\times 64$ VAE latent tensor cuts compute dramatically before final decoding.
\end{StudySolution}

来源：Lecture10b，第 194–200 个单元；Lecture10c，第 1–23 个单元（DDPM、前向边缘公式、DDIM、潜空间扩散、CFG、GLIGEN 与 Inpainting）。


\StudyHeading{1}{6. 回顾与英文自测}
\StudyTarget{CS5489-Lecture10__review}

\textbf{Q1. 简述 AE 与 VAE 的核心区别，以及为什么 VAE 更适合做隐空间插值？ / State the main differences between an AE and a VAE, and why a VAE interpolates better in latent space.}

\textbf{答 / Answer：} AE 映射出确定性隐向量 $\mathbf{z}$；VAE 输出高斯后验分布 $(\boldsymbol{\mu},\boldsymbol{\sigma})$、在损失里加 KL 散度拉近标准高斯先验，并在训练时用重参数化从后验随机采样。随机采样配合先验约束避免了隐空间出现“空洞”和过拟合单个编码，因此插值与新样本生成更平滑。 / AE predicts a single latent vector, whereas VAE predicts a Gaussian posterior, adds a KL regularizer toward $\mathcal{N}(0,I)$, and samples latent vectors during training. Sampling plus the prior prevents overfitting to isolated codes and produces a smooth latent space for interpolation.

\textbf{Q2. 为什么 GAN 无法直接计算样本的概率密度（似然），什么条件下可以计算？ / Why can't a standard GAN compute sample likelihoods, and under what condition is likelihood computable?}

\textbf{答 / Answer：} 变换 $y=f(x)$ 的密度公式需要逆映射 $f^{-1}$ 及其导数（雅可比行列式），而一般神经网络 $f$ 并不可逆。只有当神经网络完全由可逆且可微的变换复合而成时，才能精确计算似然。 / Change of variables requires $f^{-1}$ and its Jacobian determinant, which do not exist for a general neural network. The network must be composed entirely of invertible, differentiable transformations.

\textbf{Q3. GAN 实践中为何常用 $\min_G -\mathbb{E}[\log D(G(z))]$ 替代 $\min_G \mathbb{E}[\log(1-D(G(z)))]$？什么是模式崩塌？ / Why is non-saturating GAN loss preferred over minimax loss, and what is mode collapse?}

\textbf{答 / Answer：} 早期假样本容易被判别器识破（$D(G(z))\approx 0$），原始极小极大损失梯度几乎消失，而非饱和损失此时梯度大得多、训练更稳定。模式崩塌指生成器只学会生成真实分布中一个或少数几个模态来骗过判别器，丢失了分布多样性。 / Early in training when $D(G(z))\approx 0$, the minimax loss suffers vanishing gradients whereas $-\log D(G(z))$ keeps strong gradients. Mode collapse occurs when $G$ generates only one or a few modes of a multimodal target distribution (still fooling $D$, but losing diversity).

\textbf{Q4. 比较 VAE、GAN 和扩散模型三大深度生成模型的直觉差异。 / Briefly contrast VAE, GAN, and Diffusion models.}

\textbf{答 / Answer：} VAE 用单步编解码最大化变分下界（好训、有显式隐空间和近似似然，但生成偏糊）；GAN 用生成器与判别器对抗博弈（采样快、清晰，但无似然、难训易模式崩塌）；扩散模型通过多步加噪与逐步去噪（每步有明确噪声预测监督，质量高且稳定，但多步采样开销更大）。 / VAE optimizes an ELBO with a structured latent space (stable training, somewhat blurry); GAN uses adversarial training (sharp fast samples, no likelihood, unstable/mode collapse); Diffusion learns step-by-step denoising with direct per-step supervision (high fidelity and stable, slower multi-step sampling).


\StudyHeading{2}{学习要求与难点回查}
\StudyTarget{CS5489-Lecture10__learning-priorities}

\begingroup\small\setlength{\tabcolsep}{3pt}\renewcommand{\arraystretch}{1.22}
\begin{longtable}{@{}>{\raggedright\arraybackslash}p{\dimexpr0.20\linewidth-1.50000\tabcolsep\relax}>{\raggedright\arraybackslash}p{\dimexpr0.14\linewidth-1.50000\tabcolsep\relax}>{\raggedright\arraybackslash}p{\dimexpr0.22\linewidth-1.50000\tabcolsep\relax}>{\raggedright\arraybackslash}p{\dimexpr0.44\linewidth-1.50000\tabcolsep\relax}@{}}
\toprule
主题 & 学习优先级 & 理解难度与卡点 & 掌握要求／判断依据 \\
\midrule\endfirsthead
主题 & 学习优先级 & 理解难度与卡点 & 掌握要求／判断依据 \\
\midrule\endhead
\bottomrule\endfoot
AE、DAE、MAE 与 U-Net & 核心必会 & 低，易误认 AE 须对称 & 记牢 AE 五条多选判断、DAE 去噪动机与 U-Net 跳跃拼接作用 \\
重参数化与变量替换 & 核心必会 & 中，反函数可逆条件 & 默写 $z=\mu+\sigma\epsilon$，解释为何一般 NN 不能算似然 \\
VAE 原理与单类分类 & 核心必会 & 高，ELBO 两项平衡 & 熟练对比 AE/VAE 异同，能设计苹果单类检测流程 \\
GAN 博弈与模式崩塌 & 核心必会 & 中，非饱和梯度与模式崩塌 & 写对式 10.5、$D^*$，解释非饱和优势与模式崩塌 \\
扩散模型与 CFG & 核心必会 & 中，闭式 $\bar{\alpha}_t$ 与隐空间 & 写对式 10.6–10.7、CFG 公式，分清 DDPM/DDIM/Latent 扩散 \\
\end{longtable}\endgroup


\StudyHeading{1}{Historical exam map｜按考点查历史期末题源}
\StudyTarget{CS5489-Lecture10__exam-topic-index}

口径：汇总已辨识的 3 套独立期末卷（2020B、2021A、2021B）；同卷不重复计数，扩散模型为近年新增模块。

\begingroup\small\setlength{\tabcolsep}{3pt}\renewcommand{\arraystretch}{1.22}
\begin{longtable}{@{}>{\raggedright\arraybackslash}p{\dimexpr0.24\linewidth-1.33333\tabcolsep\relax}>{\raggedright\arraybackslash}p{\dimexpr0.22\linewidth-1.33333\tabcolsep\relax}>{\raggedright\arraybackslash}p{\dimexpr0.54\linewidth-1.33333\tabcolsep\relax}@{}}
\toprule
考点 & 试卷与题号 & 考查重点与提醒 \\
\midrule\endfirsthead
考点 & 试卷与题号 & 考查重点与提醒 \\
\midrule\endhead
\bottomrule\endfoot
Autoencoder 性质 & 2021A Q6 & 降维与去噪、可用预训练层；训练慢于 PCA、无监督、不必镜像对称 \\
VAE 原理、AE 对比与单类分类 & 2020B Q12；2021A Q10；2021B Q6 & AE vs VAE 异同与插值意义；仅用苹果图训 VAE 做单类检测；VAE 后验/加噪/不需共享权重 \\
GAN 机制与似然不可算 & 2020B Q6；2021B Q13 & 重参数化、非饱和损失、无似然/无配分函数、模式崩塌；变量替换需要 $f$ 可逆且可微 \\
\end{longtable}\endgroup


\StudyHeading{1}{来源覆盖索引}
\StudyTarget{CS5489-Lecture10__sources}

依据当前课件 Lecture10a（98 个单元）、Lecture10b（200 个单元）与 Lecture10c（23 个单元）。

\begingroup\small\setlength{\tabcolsep}{3pt}\renewcommand{\arraystretch}{1.22}
\begin{longtable}{@{}>{\raggedright\arraybackslash}p{\dimexpr0.70\linewidth-1.00000\tabcolsep\relax}>{\raggedright\arraybackslash}p{\dimexpr0.30\linewidth-1.00000\tabcolsep\relax}@{}}
\toprule
原课位置 & 对应讲义 \\
\midrule\endfirsthead
原课位置 & 对应讲义 \\
\midrule\endhead
\bottomrule\endfoot
Lecture10a，第 1–71 个单元 AE、权重共享、稀疏正则、DAE 向量场与概率解释 & §1 \\
Lecture10a，第 72–98 个单元 ConvAE 隐空间插值、MAE、语义分割 FCN 与 U-Net & §1 \\
Lecture10b，第 1–83 个单元 一对多生成、重参数化、变量替换密度、VAE 与 Conv-VAE & §2、§3 \\
Lecture10b，第 84–193 个单元 GAN 极小极大分析、非饱和损失、模式崩塌、CycleGAN & §4 \\
Lecture10b 第 194–200 单元；Lecture10c 第 1–23 单元 DDPM 前向／反向、DDIM、潜空间扩散、CFG、GLIGEN 与 Inpainting & §5、§6 \\
\end{longtable}\endgroup

本讲全部手算验算与自绘教学图见同目录 \StudyInline{tools/\allowbreak{}verify\_\allowbreak{}l10\_\allowbreak{}l12\_\allowbreak{}examples.\allowbreak{}py}（独立断言结果存档于 \StudyInline{l10-l12-example-results.\allowbreak{}json}）。
